\documentclass[11pt]{article}

\usepackage[margin=1in]{geometry}
\usepackage[T1]{fontenc}
\usepackage[utf8]{inputenc}
\usepackage{lmodern}
\usepackage{microtype}
\usepackage{amsmath,amssymb,amsthm}
\usepackage{mathtools}
\usepackage{xcolor}
\usepackage{hyperref}
\usepackage{booktabs}
\usepackage{enumitem}
\usepackage{listings}

\hypersetup{
  colorlinks=true,
  linkcolor=blue!50!black,
  citecolor=blue!50!black,
  urlcolor=blue!50!black
}

\lstdefinelanguage{Lean4}{
  morekeywords={
    import,namespace,end,open,section,noncomputable,abbrev,structure,inductive,
    def,theorem,lemma,where,match,with,if,then,else,by,fun,let,in,
    simp,calc,have,show,exact,rcases,cases,constructor,intro,apply
  },
  sensitive=true,
  morecomment=[l]{--},
  morecomment=[s]{/-}{-/},
  morestring=[b]"
}

\lstset{
  language=Lean4,
  basicstyle=\ttfamily\small,
  keywordstyle=\bfseries\color{blue!50!black},
  commentstyle=\itshape\color{green!40!black},
  stringstyle=\color{red!50!black},
  showstringspaces=false,
  breaklines=true,
  frame=single,
  framerule=0.3pt,
  xleftmargin=0.5em,
  xrightmargin=0.5em,
  keepspaces=true,
  columns=fullflexible
}

\title{The MQ-Calculus in Lean~4:\\
Communication-as-Measurement, Operational and Denotational Semantics}
\author{Zar \and Oru\v{z}i\footnote{``Oru\v{z}i'' denotes AI-assisted collaboration
using Claude Code (Anthropic), ChatGPT, and Codex (OpenAI).
Formalization of the MQ-calculus by Stay and Meredith~\cite{stay-meredith-mq}.}}
\date{\today}

\begin{document}
\maketitle

\begin{abstract}
We present a Lean~4 formalization of the MQ-calculus of Stay and
Meredith~\cite{stay-meredith-mq}, comprising 1{,}413 lines across 10 files
with zero sorries and zero custom axioms.  The MQ-calculus unifies
communication and quantum measurement: a COMM event produces a binary
measurement outcome governed by Born-rule probabilities.

The formalization covers: (1)~process syntax with de~Bruijn wire indices;
(2)~the COMM reduction rule with constructive both-outcomes witness;
(3)~structural congruence (parallel laws, scope extrusion via shift);
(4)~full operational reduction with contextual closure;
(5)~a backend-parametric denotational semantics with a statevector
instantiation including post-collapse normalization; and
(6)~clause-level theorem handles (\texttt{paper63\_*}) for paper auditing.

Cross-calculus bridges connect MQ measurement branching to MORK binary-fold
outcomes and place the calculus in the broader OSLF type synthesis pipeline.
\end{abstract}

\section{Why MQ-Calculus Matters}
The calculus serves three simultaneous roles.
\begin{enumerate}[leftmargin=1.5em]
  \item It stress-tests the MeTTaIL formal stack on a nontrivial stochastic process calculus.
  \item It provides concrete benchmarks for stochastic simulators beyond the stochastic $\pi$-calculus setting.
  \item It grounds the OSLF program in physically meaningful scenarios where observations (experiments) are represented structurally in the reduction semantics.
\end{enumerate}

A central methodological point is to represent interpretational choices as program structure, not as external metaphysical commitments. Distinct quantum narratives correspond to distinct valid programs with distinct typing/behavioral footprints.

\section{Paper-to-Lean Map}
The Lean formalization is organized under
\texttt{Mettapedia/Languages/ProcessCalculi/MQCalculus/}:

\begin{center}
\small
\begin{tabular}{llr}
\toprule
\textbf{File} & \textbf{Contents} & \textbf{Lines} \\
\midrule
\texttt{Syntax.lean} & Process grammar, gate specs & 120 \\
\texttt{CommRule.lean} & COMM rule, measurement branches & 95 \\
\texttt{Reduction.lean} & Operational reduction, multi-step & 172 \\
\texttt{StructuralCongruence.lean} & SC: par laws, scope extrusion & 205 \\
\texttt{Shift.lean} & De Bruijn wire shifting & 110 \\
\texttt{Backend.lean} & QState, Born probs, collapse, backend & 248 \\
\texttt{Denotational.lean} & \texttt{denoteWith}, weighted branches & 165 \\
\texttt{PaperMap.lean} & \texttt{paper63\_*} theorem handles & 130 \\
\texttt{Interoperability.lean} & MORK binary-fold bridge & 78 \\
\texttt{MQCalculus.lean} & Barrel import & 10 \\
\midrule
\textbf{Total} & \textbf{0 sorries, 0 axioms} & \textbf{1{,}413} \\
\bottomrule
\end{tabular}
\end{center}

\section{Core Syntax in Lean}
The process grammar is encoded with de Bruijn-style wire indexing and typed gate specifications.

\begin{lstlisting}[caption={Process and gate syntax (from \texttt{Syntax.lean})}]
inductive GateOp : Type where
  | H
  | X
  | Z
  | custom (name : String)
  deriving DecidableEq, Repr

structure GateSpec where
  op : GateOp
  target : Nat
  deriving DecidableEq, Repr

def gateSpecOfString : String -> GateSpec := GateSpec.ofString

inductive Process : Type where
  | MQNil  : Process
  | MQPar  : Process -> Process -> Process
  | MQNu   : Process -> Process
  | MQGate : GateSpec -> Process -> Process
  | MQOut  : Nat -> Process
  | MQIn   : Nat -> Process -> Process -> Process
  deriving DecidableEq, Repr
\end{lstlisting}

This representation keeps the surface small while supporting typed gate routing and explicit branch continuations for measurement outcomes.

\section{Communication as Measurement}
The communication rule is represented by a branch type plus a small-step relation that produces both outcome branches constructively.

\begin{lstlisting}[caption={COMM branching object and witnesses (from \texttt{CommRule.lean})}]
inductive Outcome : Type where
  | zero : Outcome
  | one  : Outcome
  deriving DecidableEq

structure MeasurementBranch where
  outcome : Outcome
  result  : Process

inductive CommReduction : Nat -> Process -> Process -> MeasurementBranch -> Prop where
  | outcome_zero (i : Nat) (p q : Process) :
      CommReduction i p q (.zero, p)
  | outcome_one (i : Nat) (p q : Process) :
      CommReduction i p q (.one, q)

theorem comm_both_outcomes (i : Nat) (p q : Process) :
    CommReduction i p q (.zero, p) /\ CommReduction i p q (.one, q) :=
  (.outcome_zero i p q, .outcome_one i p q)
\end{lstlisting}

Born-style probabilities are defined separately on vectors. The branch relation itself remains structural and proof-friendly.

\section{Operational Reduction Relation}
Full one-step reduction includes COMM firing, structural congruence closure, and contextual closure under parallel/new/gate constructors.

\begin{lstlisting}[caption={One-step and multi-step reduction (from \texttt{Reduction.lean})}]
inductive Reduces : Process -> Process -> Prop where
  | comm (i : Nat) (p q : Process) (b : MeasurementBranch) :
      CommReduction i p q b ->
      Reduces (MQPar (MQOut i) (MQIn i p q)) b.result
  | sc (p p' q' q : Process) :
      SC p p' -> Reduces p' q' -> SC q' q -> Reduces p q
  | par_l  (p p' q : Process) : Reduces p p' -> Reduces (MQPar p q) (MQPar p' q)
  | par_r  (p q q' : Process) : Reduces q q' -> Reduces (MQPar p q) (MQPar p q')
  | nu_step  (p p' : Process) : Reduces p p' -> Reduces (MQNu p) (MQNu p')
  | gate_step (g : GateSpec) (p p' : Process) :
      Reduces p p' -> Reduces (MQGate g p) (MQGate g p')

inductive MultiStep : Process -> Process -> Prop where
  | refl (p : Process) : MultiStep p p
  | step (p q r : Process) : Reduces p q -> MultiStep q r -> MultiStep p r
\end{lstlisting}

\section{Backend-Parametric Denotational Semantics}
The semantics is parameterized by a backend interface with explicit operations for gate action, fresh allocation, branch probabilities, and collapse.

\begin{lstlisting}[caption={Runtime state and backend interface (from \texttt{Backend.lean})}]
structure QState where
  n : Nat
  psi : QVec n

def rawOutcomeProb (i : Nat) (out : Outcome) (st : QState) : Real :=
  Finset.univ.sum (fun k : Fin (2 ^ st.n) =>
    if bitMatches i k.1 out then Complex.normSq (st.psi.ofLp k) else 0)

def collapseByOutcome (i : Nat) (out : Outcome) (st : QState) : QState :=
  let p := rawOutcomeProb i out st
  if p = 0 then st else
    let z : Complex := ((Real.sqrt p) : Complex)
    let psi' : QVec st.n :=
      Finset.univ.sum (fun k : Fin (2 ^ st.n) =>
        EuclideanSpace.single k
          (if bitMatches i k.1 out then st.psi.ofLp k / z else 0))
    (st.n, psi')

structure MQSemanticsBackend where
  branchProb : Nat -> Outcome -> QState -> Real
  branchProb_nonneg : forall i out st, 0 <= branchProb i out st
  branchProb_sum_one : forall i st,
    branchProb i .zero st + branchProb i .one st = 1
  applyGate : GateSpec -> QState -> QState
  allocFresh : QState -> QState
  collapse : Nat -> Outcome -> QState -> QState
\end{lstlisting}

The current \texttt{statevectorBackend} instantiation gives executable semantics while preserving theorem-level invariants such as nonnegativity, normalization splits, and collapse behavior.

\section{Executable Denotation and Measurement Selection}
The denotation function is structurally recursive on process syntax and resolves \texttt{MQIn} by comparing branch probabilities.

\begin{lstlisting}[caption={Denotational clauses and weighted measurement view (from \texttt{Denotational.lean})}]
def denoteWith (backend : MQSemanticsBackend) : Process -> QState -> Option QState
  | .MQNil,      st => some st
  | .MQOut _,    st => some st
  | .MQGate s p, st => denoteWith backend p (backend.applyGate s st)
  | .MQNu p,     st => denoteWith backend p (backend.allocFresh st)
  | .MQPar p q,  st => (denoteWith backend p st).bind (denoteWith backend q)
  | .MQIn i p q, st =>
      if backend.branchProb i .zero st >= backend.branchProb i .one st then
        denoteWith backend p (backend.collapse i .zero st)
      else
        denoteWith backend q (backend.collapse i .one st)

structure WeightedBranch where
  prob_zero    : Real
  prob_one     : Real
  state_zero   : QState
  state_one    : QState
  prob_zero_nn : 0 <= prob_zero
  prob_one_nn  : 0 <= prob_one
  prob_sum_le  : prob_zero + prob_one <= 1
\end{lstlisting}

This gives two aligned views: deterministic branch selection for executable denotation and explicit two-branch packaging for proof/audit work.

\section{Clause-Level Theorem Map}
The theorem index module makes paper auditing straightforward.

\begin{lstlisting}[caption={Paper-clause theorem handles (from \texttt{PaperMap.lean})}]
theorem paper63_nil ...
theorem paper63_out ...
theorem paper63_gate ...
theorem paper63_new ...
theorem paper63_par ...
theorem paper63_in ...
theorem paper63_measurement_probs_sum_one ...
theorem paper63_measurement_state_zero ...
theorem paper63_measurement_state_one ...
theorem paper63_statevector_collapse_norm_one_of_raw_pos ...
theorem paper63_comm_denote_eq_measurement_selection ...
\end{lstlisting}

\noindent
A key normalization result used in review is:
\begin{lstlisting}[caption={Post-collapse normalization (from \texttt{Backend.lean})}]
theorem collapseByOutcome_norm_eq_one_of_raw_pos
    (i : Nat) (out : Outcome) (st : QState)
    (hraw : 0 < rawOutcomeProb i out st) :
    norm ((collapseByOutcome i out st).psi) = 1
\end{lstlisting}

\section{Cross-Calculus Bridges}

\subsection{MQ and MORK}
MQ measurement branching is linked to MORK binary-fold branching via an explicit
bridge theorem:

\begin{lstlisting}[caption={MQ--MORK bridge (from \texttt{Interoperability.lean})}]
theorem comm_nondeterminism_iff_mork_binary (i : Nat) (p q : Process) :
    (CommReduction i p q (.zero, p) /\ CommReduction i p q (.one, q)) <->
    forall (fold : FoldStep) (_hb : fold.isBinary),
      (exists (f0 : FoldStep) (_ : f0.isBinary),
        FoldPicksSubResult f0 hbin .subResult0) /\
      (exists (f1 : FoldStep) (_ : f1.isBinary),
        FoldPicksSubResult f1 hbin .subResult1)
\end{lstlisting}

This ties communication-level nondeterminism in MQ directly into the MORK
three-phase execution model.

\subsection{Shared Process Calculus Infrastructure}

The Mettapedia formalization contains three process calculi---$\pi$-calculus
(8{,}768 lines), $\rho$-calculus (4{,}848 lines), and MQ-calculus---that
share $\sim$40\% of their reduction/congruence infrastructure.  A shared
library in \texttt{ProcessCalculi/Common/} eliminates this duplication:

\begin{center}
\small
\begin{tabular}{llr}
\toprule
\textbf{File} & \textbf{Contents} & \textbf{Lines} \\
\midrule
\texttt{Common/Star.lean} & \texttt{RTClosure} (Type) + \texttt{RTClosureProp} (Prop) closures & 140 \\
 & Generic lifters: \texttt{liftUnary}, \texttt{liftBinLeft}, \texttt{liftBinRight} & \\
\texttt{Common/Congruence.lean} & \texttt{HasSC} typeclass, \texttt{reduces\_modulo\_sc} & 65 \\
\texttt{Common/ProcessAlgebra.lean} & \texttt{HasPar}, \texttt{HasNu}, \texttt{HasNil} typeclasses & 30 \\
\bottomrule
\end{tabular}
\end{center}

Each calculus registers instances:
\begin{itemize}[leftmargin=1.5em]
  \item \textbf{$\pi$-calculus}: \texttt{HasPar}, \texttt{HasNil}, \texttt{HasSC}
    (wrapping Type-valued SC in \texttt{Nonempty}).
  \item \textbf{$\rho$-calculus}: \texttt{HasSC} (Prop-valued SC directly).
    The $\rho$-calculus also \emph{gained} multi-step parallel congruence lifting
    (\texttt{ReducesStar.par\_head}, \texttt{ReducesStar.par\_any\_pos}) that was
    previously missing.
  \item \textbf{MQ-calculus}: \texttt{HasPar}, \texttt{HasNil}, \texttt{HasNu},
    \texttt{HasSC}.
\end{itemize}

All three calculi share:
\begin{itemize}[leftmargin=1.5em]
  \item \textbf{Structural congruence}: equivalence relation with parallel
    commutativity/associativity and scope extrusion, unified via \texttt{HasSC}.
  \item \textbf{COMM rule shape}: sender + receiver $\to$ result,
    differing only in payload consumption (substitution in $\pi$/$\rho$,
    measurement in MQ).
\end{itemize}

\subsection{Spice Calculus and Present Moment}

The Spice calculus (formalized in \texttt{RhoCalculus/SpiceRule.lean})
extends $\rho$-calculus COMM with $n$-step lookahead: agents receive
the set of states reachable in $\leq n$ steps before committing to an
interaction.  When $n = 0$ this recovers standard $\rho$-calculus COMM;
the ``present moment'' ($n = 1$) gives agents awareness of their immediate
interaction possibilities.

The new \texttt{ReducesN.par\_head} and \texttt{ReducesN.par\_any\_pos}
theorems (n-step parallel congruence lifting) enable compositional
reasoning about Spice agents: an agent's future states lift cleanly through
parallel contexts (\texttt{futureStates\_par\_head},
\texttt{presentMoment\_par\_head}, \texttt{reachableStates\_par\_head}).
This is essential for the Spice rule, where agents must evaluate their
environment's $n$-step evolution.

\section{Wigner-Style Program Distinctions}
A practical interpretational takeaway is representational: one can encode a fully unitary evolution program and, separately, a staged-communication program where measurement-triggering communications occur between stages. In this framework they are different programs with different transition structures. The formal machinery does not force them to collapse into a single execution narrative.

\section{OSLF Context and AGI-Relevant Direction}
The Leifer--Milner--Sewell context-generation perspective provides the right conceptual anchor: contexts act as experiments. In MQ settings this dovetails with communication-as-measurement and gives a principled route from rewrite semantics to observational structure.

For AGI-motivated scientific discovery workflows, the important systems consequence is that physical theories can be encoded as rewrite theories, transformed, and evaluated in a semantics-aware loop. Extending from MQ-calculus to additional physics rewrite systems (e.g., GR-oriented systems) becomes a matter of language-instance construction and bridge theorems, not ad hoc simulator rewrites.

\section{Next Formalization Targets}
\begin{enumerate}[leftmargin=1.5em]
  \item \textbf{Multi-qubit gate semantics} beyond single-wire gate action.
  \item \textbf{Stronger collapse theory}: preservation under composition,
    idempotence of repeated measurement on the same wire.
  \item \textbf{OSLF integration}: MQ process experiments as OSLF language
    instances with $\Diamond \dashv \Box$ Galois connection, paralleling the
    PeTTa and Governance OSLF instances.
  \item \textbf{Spice--MQ bridge}: connect the $\rho$-calculus Spice
    $n$-step lookahead to MQ measurement prediction (agents observing
    future measurement outcomes).
\end{enumerate}

\begin{thebibliography}{9}
\bibitem{stay-meredith-mq}
M.~Stay and L.~G.~Meredith,
\emph{The MQ-Calculus: A Calculus for Quantum Processes Involving Intermittent
  Measurement},
2026 working draft.

\bibitem{meredith-present-moment}
L.~G.~Meredith,
\emph{How the Agents Got Their Present Moment},
2026.

\bibitem{meredith-radestock-2005}
L.~G.~Meredith and M.~Radestock,
\emph{A Reflective Higher-Order Calculus},
\textit{Electronic Notes in Theoretical Computer Science}, 141(5), 2005.

\bibitem{gay-nagarajan-cqp}
S.~J.~Gay and R.~Nagarajan,
\emph{Communicating Quantum Processes},
POPL, 2005.

\bibitem{leifer-milner}
J.~Leifer and R.~Milner,
\emph{Deriving Bisimulation Congruences for Reactive Systems},
CONCUR, 2000.
\end{thebibliography}

\end{document}
